\origpage{123}{GRAPH MINORS. XV. GIANT STEPS}
\hspace*{12pt}(iii) \textit{$\theta \le ord(\mathscr{T}')$, and $(A \cap H', B \cap H') \in \mathscr{T}'$ for every $(A, B) \in \mathscr{T}^*$ of order $< \theta$.}

(1)(ii) is clear, and (1)(iii) is proved in the same way as in the proof of (3.3). It remains to show (1)(i). Let $\Sigma_1 = \Sigma - \Lambda$, $\Sigma_2 = \bar{\Lambda}$ and $X = \{v_1, \dots, v_{2k}\}$. We shall verify the hypotheses of (3.2). Now (3.2)(i) is trivial since $\Sigma_1 \cap \Sigma_2$ is connected; and (3.2)(ii) is proved as in the proof of (3.3). For (3.2)(iii), let $F \subseteq \Sigma_2$ be an $O$-arc with $F \cap U(H') \subseteq V(H')$ and $|F \cap V(H')| \le 2$; then the closed disc in $\Sigma_2$ bounded by $F$ is a dial for $F$ because $H' \cap \Sigma_2$ is the union of $H \cap bd(\Lambda)$ and the disjoint paths $P_1, \dots, P_{k-1}$. For (3.2)(iv), let $u, v, F, P$ be as in (3.2)(iv). If $F \subseteq r_1 \cup \{u, v\}$ or $F \subseteq r_2 \cup \{u, v\}$, then (3.2)(iv) is satisfied, and so we may assume that $F \subseteq r_3 \cup \{u, v\}$ for some region $r_3 \ne r_1, r_2$ of $H'$ with $r_3 \subseteq \Lambda$. Let the four vertices in $X$ incident with $r_3$ be $v_i, v_{i+1}, v_{2k-i-1}, v_{2k-i}$, where $1 \le i \le k - 2$. Now there is a path of $H \cap \Sigma$ between $u$ and $v$, passing through at most two vertices in $X$; and so $u, v$ both lie in the same component of $H \cap bd(r_3)$ (because of the vertices $v_k, v_{2k} \in X$) and again (3.2)(iv) is satisfied. \fix{We give the details. First, for $1 \le j \le k-1$ we show that $U(P_j) \subseteq \Lambda \cup \{v_j, v_{2k-j}\}$. Let $e$ be an edge of $P_j$. Since $U(P_j) \cap U(H) = \{v_j, v_{2k-j}\}$ and these are vertices, $e \cap U(H) = \emptyset$, and so $e \cap bd(\Lambda) = \emptyset$. Since $e$ is connected, either $e \subseteq \Lambda$ or $e \subseteq \Sigma - \bar{\Lambda}$; and the second is impossible, since then $e$ would be an edge of $H' \cap (\Sigma - \Lambda) = H$. Thus every edge of $P_j$ is a subset of $\Lambda$, and so every vertex of $P_j$ is in $\bar{\Lambda}$. Every vertex of $P_j$ other than $v_j, v_{2k-j}$ is not in $U(H)$, and hence not in $bd(\Lambda)$, and so it is in $\Lambda$. This proves that $U(P_j) \subseteq \Lambda \cup \{v_j, v_{2k-j}\}$. Since $U(H) \cap \Lambda = \emptyset$, it follows that $\Sigma_2 \cap U(H') = bd(\Lambda) \cup U(P_1) \cup \dots \cup U(P_{k-1})$, and that $H' \cap bd(\Lambda) = H \cap bd(\Lambda)$.

Now let $u, v, F, P$ be as in (3.2)(iv). We show first that for $1 \le j \le k-1$, $P$ does not contain both $v_j$ and $v_{2k-j}$. The circuit $H \cap bd(\Lambda)$ is the union of two paths with ends $v_j, v_{2k-j}$, and since $v_1, \dots, v_{2k}$ are in order on this circuit, one of these paths contains $v_k$ and the other contains $v_{2k}$. If $v_j, v_{2k-j} \in V(P)$, then since $P$ is a path of this circuit, $P$ includes one of these two paths, and so $|V(P) \cap X| \ge 3$, a contradiction.

Since $F$ meets $U(H')$ only in $u$ and $v$, the set $U(P) \cup F$ is an $O$-arc in the closed disc $\Sigma_2$, and so it bounds a closed disc $D \subseteq \Sigma_2$. Let $D^0 = D - (U(P) \cup F)$. Then $D^0$ is an open disc, and $D^0 \subseteq \Lambda$, since $D^0$ is open in $\Sigma$ and $D^0 \subseteq \bar{\Lambda}$. We claim that $D^0 \cap U(H') = \emptyset$. For suppose that $x \in D^0 \cap U(H')$. Since $x \in \Lambda$, the first paragraph implies that $x \in U(P_j) - \{v_j, v_{2k-j}\}$ for some $j$. This set is connected, and it is disjoint from $U(P) \cup F$, since $U(P) \cup (F \cap U(H')) \subseteq bd(\Lambda)$. Since it meets $D^0$ and $\Sigma - D$ is open, it is a subset of $D^0$, and so $v_j, v_{2k-j} \in D$. But $v_j, v_{2k-j} \in bd(\Lambda)$, and so they are not in $D^0$. Hence they belong to $U(P) \cup F$, and therefore to $U(P)$, since $F \cap U(H') = \{u, v\} \subseteq U(P)$. Thus $v_j, v_{2k-j} \in V(P)$, contrary to what we showed above. This proves that $D^0 \cap U(H') = \emptyset$, and so (3.2)(iv) is satisfied.}\fixnote{fix:3.4-iv-step-added}{The paper does not prove that the vertices of $X$ incident with $r_3$ are $v_i, v_{i+1}, v_{2k-i-1}, v_{2k-i}$, and it does not say why (3.2)(iv) holds when $F \subseteq r_1 \cup \{u, v\}$ or $F \subseteq r_2 \cup \{u, v\}$. The argument above is the one used in the formalization.}

By (3.2), $H'$ is rigid. Hence $H', \mathscr{T}'$ is a $\Sigma$-span of order $\theta$. Let $d'$ be the metric of $\mathscr{T}'$, and let $e_1, e_2$ be edges of $P_1, P_{k-1}$ respectively. Now $ord(\mathscr{T}') \ge 2k > 2(k - 1)$, and so by [4, Theorem (6.1)], $d'(e_1, e_2) \ge 2k - 2$. \fix{We give the details. By the first paragraph of the details for (3.2)(iv) above, $U(P_1), \dots, U(P_{k-1})$ are mutually disjoint lines in the closed disc $\bar{\Lambda}$, and $U(P_j) \cap bd(\Lambda) = \{v_j, v_{2k-j}\}$ for $1 \le j \le k-1$. Let $D_0$ be the closed unit disc in the plane, and for $1 \le j \le k-1$ let $C_j$ be the vertical chord of $D_0$ with ends $(\cos(\pi j/k), \pm \sin(\pi j/k))$. Thus $C_1, \dots, C_{k-1}$ are disjoint, and they occur in this order from right to left.

We claim that there is a homeomorphism $\phi \colon \bar{\Lambda} \to D_0$ that maps $v_i$ to $(\cos(\pi i/k), \sin(\pi i/k))$ for $1 \le i \le 2k$, and maps $U(P_j)$ onto $C_j$ for $1 \le j \le k-1$. Since $v_1, \dots, v_{2k}$ are in order on the $O$-arc $bd(\Lambda)$, there is a homeomorphism $\phi$ with the first property. Suppose that $\phi(U(P_j)) = C_j$ for $1 \le j < n$, where $1 \le n \le k-1$. Let $D_1 = D_0$ if $n = 1$, and otherwise let $D_1$ be the part of $D_0$ on the left of $C_{n-1}$, including $C_{n-1}$. The line $L = \phi(U(P_n))$ is disjoint from $C_{n-1}$, and its ends are on the left of $C_{n-1}$, so $L \subseteq D_1$. Both $L$ and $C_n$ are lines in the closed disc $D_1$ with the same ends, meeting $bd(D_1)$ only in their ends. Each of them divides $D_1$ into two closed discs, bounded by the line and by the two arcs of $bd(D_1)$ between its ends. Since a homeomorphism between the boundaries of two closed discs extends to a homeomorphism between the discs, there is a homeomorphism of $D_1$ onto itself that fixes $bd(D_1)$ and maps $L$ onto $C_n$. We extend it by the identity on $D_0 - D_1$, and compose $\phi$ with the result. The new $\phi$ still maps $U(P_j)$ onto $C_j$ for $j < n$, since $C_j \subseteq (D_0 - D_1) \cup bd(D_1)$, and now also for $j = n$. The claim follows by induction on $n$.

The image under $\phi^{-1}$ of the part of $D_0 - bd(D_0)$ on the left of $C_{k-1}$ is connected and open, it is disjoint from $U(H')$, and its boundary is included in $U(H')$. Hence it is a region of $H'$ in $\Lambda$. It is incident with $v_k$, and so it is $r_1$, and it is incident with every edge of $P_{k-1}$. Similarly $r_2$ is the image of the part of $D_0 - bd(D_0)$ on the right of $C_1$, and it is incident with $v_{2k}$ and with every edge of $P_1$.

Now let $\psi$ be a homeomorphism of the plane with $\psi(D_0) \subseteq D_0$ that fixes every point of $C_1 \cup \dots \cup C_{k-1}$ and maps every other point of $bd(D_0)$ into $D_0 - bd(D_0)$. (For instance, let $c = \cos(\pi/k)$, let $h(x) = x$ if $|x| \le c$ and $h(x) = (x + c)/2$ if $x > c$ and $h(x) = (x - c)/2$ if $x < -c$, let $s(x) = 1 - \min(1/2, \prod_{j=1}^{k-1} (x - \cos(\pi j/k))^2)$, and let $\psi(x, y) = (h(x), s(x)y)$.) Let $\Delta = \phi^{-1}(\psi(D_0))$ and $F = bd(\Delta)$. Then $\Delta \subseteq \bar{\Lambda}$ is a closed disc, $\Delta \cap bd(\Lambda)$ is the set $Y$ of the $2k-2$ ends of $P_1, \dots, P_{k-1}$, and every other point of $U(P_1) \cup \dots \cup U(P_{k-1})$ is in $\Delta - F$. (For $\psi(D_0) \cap bd(D_0)$ is the set of ends of $C_1, \dots, C_{k-1}$, and $\psi$ fixes each $C_j$ and maps $D_0 - bd(D_0)$ onto $\psi(D_0) - bd(\psi(D_0))$.) Since $U(H) \cap \bar{\Lambda} \subseteq bd(\Lambda)$, it follows that $H' \cap \Delta = P_1 \cup \dots \cup P_{k-1}$, that $H' \cap \overline{\Sigma - \Delta} = H$, and that $F \cap U(H') = Y$, so $F$ is $H'$-normal. The members of $Y$ occur on $F$ in the order $v_1, \dots, v_{k-1}, v_{k+1}, \dots, v_{2k-1}$, since $\psi$ fixes them and maps $bd(D_0)$ onto $bd(\psi(D_0))$; and $P_j$ has ends the $j$th and the $(2k-1-j)$th of them. Moreover $H'$ and $H$ are 2-cell, since they are rigid; $\mathscr{T}$ is a respectful tangle in $H$ of order $\theta \ge 2k > 2k-2$, where $\theta \ge 2k$ since $X$ is free with respect to $\mathscr{T}$; $\mathscr{T}'$ is the tangle in $H'$ induced by $\mathscr{T}$; and $Y \subseteq X$ is free with respect to $\mathscr{T}$, as in the proof of (2) in the proof of (3.3). Since $2k - 2 \ge 4$, [4, Theorem (6.1)] applies with $\Gamma = H'$, with $2k-2$ in place of $\theta$, and with $\Gamma_1 = H$ and $\mathscr{T}_1 = \mathscr{T}$, and it gives $d'(e_1, e_2) \ge 2k-2$.}\fixnote{fix:3.4-6.1-step-added}{The paper does not give the disc to which [4, Theorem (6.1)] is applied, and does not check the hypotheses of that result. The argument above is the one used in the formalization.} Hence $d'(r_1, r_2) \ge 2k - 6$, since $d'(r_1, e_{\delm{1}\fix{2}}) \le 2$ and $d'(r_2, e_{\delm{2}\fix{1}}) \le 2$; and $d'(v_{\delm{1}\fix{k}}, v_{\delm{2}\fix{2k}}) \ge 2k - 8$, since $d'(v_{\delm{1}\fix{k}}, e_{\delm{1}\fix{2}}) \le 3$ and $d'(v_{\delm{2}\fix{2k}}, e_{\delm{2}\fix{1}}) \le 3$.\fixnote{fix:3.4-proof-indices}{Here $r_1$ is incident with the edges of $P_{k-1}$, among them $e_2$, and $r_2$ with the edges of $P_1$, among them $e_1$. The statement of (3.4) bounds $d'(v_k, v_{2k})$.} $\square$

Our third application of (3.2) is the following, which will be needed in a future paper.

\stmt
(3.5) \textit{Let $\Sigma$ be a surface and let $\phi \ge 3$ be an integer. Let $\mathscr{T}^*$ be a tangle in a graph $G$, and let $H, \eta, \mathscr{T}$ be a $\Sigma$-span. Let $\Lambda_1, \Lambda_2 \subseteq \Sigma$ be distinct regions of $H$ so that $H \cap bd(\Lambda_i)$ is a circuit $C_i$ ($i = 1, 2$), and no region of $H$ is incident with a vertex in $C_1$ and with a vertex in $C_2$\fix{, and every $H$-normal line $F \subseteq \Sigma - (\Lambda_1 \cup \Lambda_2)$ with one end in $bd(\Lambda_1)$ and the other in $bd(\Lambda_2)$ satisfies $|F \cap V(H)| \ge \phi$}.\fixnote{fix:3.5}{The proof uses Graph Minors XII, (5.1), which needs the added hypothesis. Without it, (3.5) is false. In the sphere, take $C_1$ and $C_2$ joined by a line that meets only 4 vertices of $H$, and take $\phi \ge 5$. After the handle is added, this line and a line along $P_1$ form an $O$-arc that does not separate the torus but meets $H'$ in fewer than $\phi$ vertices, so no tangle of order $\phi$ in $H'$ is respectful. With the hypothesis, (3.5) holds. At its only use, in Claim (4) in the proof of 8.1 of Graph Minors XVI, page 72, $\Lambda_1$ and $\Lambda_2$ come from cleared zones around two vertices that are far apart in the metric of the tangle, and this gives the hypothesis.} Let $a_1, \dots, a_\phi$ be distinct vertices of $C_1$ in order, and let $b_1, \dots, b_\phi$ be distinct vertices of $C_2$ in order (under some orientations of $C_1$ and $C_2$), such that $\{a_1, \dots, a_\phi\}$ and $\{b_1, \dots, b_\phi\}$ are free with respect to $\mathscr{T}$. For $1 \le i \le \phi$ let $P_i$ be an $\eta(H)$-path with ends $\eta(a_i), \eta(b_i)$, such that $P_1, \dots, P_\phi$ are mutually disjoint. Then there is a $\Sigma'$-span of order $\phi$, where $\Sigma'$ is a surface obtained by adding a handle to $\Sigma$.}

\textit{Proof.} We may assume that $\eta$ is the identity. Let $\Sigma_1 = \Sigma - (\Lambda_1 \cup \Lambda_2)$. Let $\Sigma_2$ be homeomorphic to a closed cylinder, with boundary $U(C_1) \cup U(C_2)$, such that $\Sigma' = \Sigma_1 \cup \Sigma_2$ is a surface obtained from $\Sigma$ by adding a handle, and there are $\phi$ disjoint lines in $\Sigma_2$ with ends $a_i, b_i$ ($1 \le i \le \phi$). We may therefore assume that $P_1, \dots, P_\phi$ are drawings in $\Sigma_2$ and hence in $\Sigma'$. Let $H' = H \cup P_1 \cup \dots \cup P_\phi$, let $\mathscr{T}''$ be the tangle in $H'$ of order $\theta$ induced by $\mathscr{T}$, and let $\mathscr{T}'$ be the $\phi$-truncation of $\mathscr{T}''$. We claim that
